\documentclass[11pt]{article}
\usepackage{shared/luxcover}
\usepackage{amsmath,amssymb,amsthm}
\usepackage{graphicx}
\usepackage{booktabs}
\usepackage{hyperref}
\usepackage{listings}
\usepackage{xcolor}
\usepackage{tikz}
\usepackage{siunitx}
\usepackage{multirow}
\usepackage{algorithm}
\usepackage{algpseudocode}
\usetikzlibrary{arrows.meta,positioning,calc,shapes.geometric,fit}

\sisetup{group-separator={,}}

\definecolor{luxblue}{RGB}{30,64,175}
\definecolor{luxgold}{RGB}{234,179,8}
\definecolor{dagteal}{RGB}{13,148,136}
\definecolor{linearred}{RGB}{220,38,38}

\newtheorem{theorem}{Theorem}
\newtheorem{lemma}[theorem]{Lemma}
\newtheorem{corollary}[theorem]{Corollary}
\theoremstyle{definition}
\newtheorem{definition}[theorem]{Definition}

\title{DAG-EVM and Conflict-Graph Consensus\\for the Lux Multi-Chain}
\author{
Zach Kelling\thanks{Corresponding author: zach@lux.network}\\
\textit{Lux Industries, Inc.}\\
\texttt{research@lux.network}
}
\date{March 2026}

\begin{document}
\luxcoverpage

\begin{abstract}
We present a unified DAG consensus model for the Lux multi-chain that lifts
Block-STM-style optimistic concurrency from a per-block scheduler into a
first-class finality primitive. Each block-producing chain (C-Chain EVM,
D-Chain DEX, Z-Chain ZK-UTXO, A-Chain AI, O-Chain Oracle, R-Chain Relay,
S-Chain ServiceNode) emits transaction batches as vertices in a native
directed acyclic graph, with conflict sets derived from the chain's natural
state model: EVM read/write sets, UTXO nullifiers, DEX order tuples,
job identifiers, feed rounds, destination-chain nonces, or service epochs.
Antichains are finalised in parallel under the Quasar post-quantum
certificate scheme; sequential worldview is preserved by a serialisability
theorem. We formalise the safety and no-double-spend properties in Lean~4
and TLA$^+$, prove commutativity of non-conflicting vertices, and show that
GPU-native conflict-graph reduction plus GPU ecrecover/keccak yields a
$32\times$ practical speedup on the dominant critical-path operations. The
result is a DAG-EVM that preserves bytecode-level compatibility with
Ethereum, a DAG-Z-Chain that preserves ZK-UTXO double-spend security, and a
framework that cleanly contrasts Lux against contemporary parallel and
DAG blockchains (Aptos, Sui, Solana, Monad, Fuel).
\end{abstract}

\tableofcontents
\newpage

%% ============================================================================
\section{Introduction}
%% ============================================================================

Parallel blockchain execution splits into two orthogonal problem families:
\emph{intra-block} parallelism (execute the transactions of a linear block
in parallel) and \emph{inter-block} parallelism (finalise non-conflicting
blocks in parallel, a partial-order consensus). The industry has focused
almost exclusively on the first: Aptos' Block-STM~\cite{blockstm}, Monad's
pipelined execution~\cite{monad}, Sei's parallel EVM, and Polygon Miden's
optimistic execution all run in linear-chain consensus. DAG consensus, when
present (Sui's Mysticeti~\cite{mysticeti}, Aleph Zero's AlephBFT,
Avalanche's Snowman$^{++}$), typically does not thread the DAG into the
state-machine layer: blocks are ordered by the DAG and then executed
sequentially, so the DAG exists above the execution engine rather than
inside it.

Lux takes the third path: \emph{the execution engine's conflict graph
\textbf{is} the consensus DAG}. Each VM emits vertices whose parents are
determined by the state-level dependency relation of that VM, and the DAG
engine (nebula) finalises antichains that by construction contain no
mutually-conflicting vertices. The downstream property is
\textbf{serialisability by construction}: any topological order of the
accepted DAG yields identical state, and no two finalised siblings can
double-spend, re-write the same EVM slot, or re-submit the same nullifier.

This paper presents:

\begin{itemize}
  \item A uniform definition of the per-chain conflict rule
    (\S\ref{sec:conflict});
  \item The DAG-EVM construction and its Block-STM lineage
    (\S\ref{sec:dagevm});
  \item The DAG-Z-Chain nullifier-conflict construction
    (\S\ref{sec:dagzchain});
  \item Lean~4 formal proofs of safety, commutativity, and
    serialisability (\S\ref{sec:lean});
  \item TLA$^+$ specifications and model-checked invariants
    (\S\ref{sec:tla});
  \item A principled comparison with contemporary parallel and DAG
    blockchains (\S\ref{sec:contemporaries});
  \item A discussion of GPU-native conflict detection and its measured
    speedup (\S\ref{sec:gpu}).
\end{itemize}

%% ============================================================================
\section{Chain-Level Conflict Rules}
\label{sec:conflict}
%% ============================================================================

Every VM in the Lux fleet exposes a single predicate,
$\mathsf{conflicts}(v_1, v_2) \in \{\mathsf{true},\mathsf{false}\}$,
over pairs of vertices. The predicate must be \emph{symmetric} and
\emph{computable from the vertex batch alone}. Table~\ref{tab:conflict}
summarises the per-chain rule.

\begin{table}[h]
\centering
\small
\begin{tabular}{lllll}
\toprule
\textbf{Chain} & \textbf{Symbol} & \textbf{Engine} & \textbf{Conflict key set} & \textbf{Source} \\
\midrule
C-Chain (EVM)         & C & DAG  & $R(v) \cup W(v) \subseteq \mathsf{StorageKey}$ & Block-STM \\
D-Chain (DEX)         & D & DAG  & $\{(base,quote,side,id)\} \subseteq \mathsf{OrderKey}$ & orderbook \\
X-Chain (UTXO)        & X & DAG  & $\{\mathsf{txo}_i\} \subseteq \mathsf{UTXO}$ & UTXO inputs \\
Z-Chain (ZK-UTXO)     & Z & DAG  & $\{\mathsf{null}_i\} \subseteq \mathsf{Nullifier}$ & nullifier reveal \\
Q-Chain (PQ)          & Q & DAG+Quasar & $\{\mathsf{stamp}_i\}$ & PQ attestation \\
A-Chain (AI)          & A & DAG  & $\{\mathsf{jobID}_i\}$ & job graph \\
O-Chain (Oracle)      & O & DAG  & $\{(\mathsf{feed},\mathsf{round})\}$ & feed rounds \\
R-Chain (Relay)       & R & DAG  & $\{(\mathsf{dstChain},\mathsf{nonce})\}$ & relay nonce \\
S-Chain (ServiceNode) & S & DAG  & $\{(\mathsf{node},\mathsf{epoch})\}$ & service epoch \\
\midrule
P-Chain (Platform)    & P & Linear & ordered         & validator set \\
B-Chain (Bridge)      & B & Linear & ordered         & ceremony steps \\
T-Chain (Threshold)   & T & Linear & ordered         & FROST rounds \\
M-Chain (MPC)         & M & Sharded-linear & per-wallet linear & CGGMP21 rounds \\
K-Chain (Keys)        & K & Linear & ordered         & key rotation \\
I-Chain (Identity)    & I & Linear & ordered         & identity updates \\
\bottomrule
\end{tabular}
\caption{Per-chain conflict rule. The first block operates under the DAG
engine with native parallel finality. The second block retains linear
ordering where the state machine requires total order (validator registry,
cryptographic ceremonies, identity updates).}
\label{tab:conflict}
\end{table}

We say a pair of vertices \emph{commutes} iff it does not conflict.

\begin{definition}[Conflict DAG]
A \emph{conflict DAG} $D = (V, E, \mathsf{conflicts})$ is a directed acyclic
graph whose vertices are VM-specific transaction batches, whose edges are
causal parent relationships, and whose conflict predicate satisfies the
\emph{Siblings-Commute invariant}: for every pair of vertices $v_1, v_2$
with $v_1 \not\preceq v_2$ and $v_2 \not\preceq v_1$ in the DAG partial
order, $\mathsf{conflicts}(v_1, v_2)$ is false.
\end{definition}

The Siblings-Commute invariant is enforced constructively by the VM's
$\mathsf{BuildVertex}$: when forming a new vertex the builder scans the
current \emph{frontier} (set of rank-maximal accepted vertices) and refuses
any batch whose conflict key set overlaps a frontier sibling. This is an
$O(|\text{frontier}| \cdot |\text{keys}|)$ bitmap intersection; on GPU it
is a single threadgroup-reduce kernel.

%% ============================================================================
\section{DAG-EVM}
\label{sec:dagevm}
%% ============================================================================

\subsection{Lineage: Block-STM Lifted}

Block-STM~\cite{blockstm} (Aptos Labs, 2022) is an optimistic concurrency
control (OCC) algorithm that speculatively executes a linear block's
transactions in parallel on separate copies of state, tracks per-transaction
read and write sets, and re-executes any transaction whose reads were
invalidated by an earlier-committed transaction's writes.

Block-STM's correctness reduces to two properties:
\begin{enumerate}
  \item Every execution is \emph{conflict-equivalent} to the sequential
    execution of the same block in the block's canonical order.
  \item Every committed write is visible to every strictly-later transaction
    that reads the same storage location.
\end{enumerate}

These properties are purely local to the block's set of transactions; they
make no reference to the fact that blocks are arranged in a linear chain.
This is the opening we exploit.

\subsection{Construction}

Let $T$ be a sequence of pending EVM transactions. The DAG-EVM builder
performs:

\begin{algorithm}[h]
\caption{BuildVertex (DAG-EVM)}
\begin{algorithmic}[1]
\State Drain up to $N$ txs from the mempool into $T$
\State Run Block-STM speculative exec on $T$ over the current state snapshot
\State Emit $(R, W) \gets$ union of per-tx read/write sets
\State $P \gets$ minimal antichain of frontier vertices whose $W$ covers $R$
\State Let $v.\mathsf{parents} \gets P$
\State $v.\mathsf{id} \gets \mathsf{keccak256}(\mathsf{txHashes} \parallel R \parallel W)$
\State Emit $v$ to the DAG engine
\end{algorithmic}
\end{algorithm}

The DAG engine (nebula) votes on acceptance. Once an antichain at rank $r$
is fully populated with accepted vertices, Quasar stamps it with a
post-quantum triple-signature certificate and $r$ becomes final. The EVM
state is updated by applying the vertices of each finalised antichain in
any valid topological order.

\subsection{Correctness Sketch}

\begin{theorem}[Non-Conflict Commutativity, Lean: \texttt{nonConflict\_commute}]
\label{thm:commute}
For any two EVM vertices $v_1, v_2$ with $\mathsf{conflicts}(v_1, v_2)
= \mathsf{false}$ and any state $\sigma$:
\[
v_2.\mathsf{apply}(v_1.\mathsf{apply}(\sigma)) = v_1.\mathsf{apply}(v_2.\mathsf{apply}(\sigma)).
\]
\end{theorem}

\begin{proof}[Proof sketch]
By case analysis on membership of each storage key $k$ in the write sets
of $v_1$ and $v_2$. The non-conflict hypothesis excludes write-write
overlap; the footprint lemmas (\texttt{writesOnly}, \texttt{readsOnly}) in
the Lean development handle the remaining four cases. Full proof in
\texttt{Consensus/DAGEVM.lean:60--135}.
\end{proof}

\begin{theorem}[Topological Order Invariance, Lean: \texttt{topo\_equivalence}]
\label{thm:topo}
Given the Siblings-Commute invariant, any two topological orderings $L_1,
L_2$ of the same accepted DAG $D$ satisfy
$\mathsf{applyList}(L_1, \sigma) = \mathsf{applyList}(L_2, \sigma)$
for every initial state $\sigma$.
\end{theorem}

\begin{theorem}[No Double Spend, Lean: \texttt{no\_double\_spend}]
If $v, w$ are both accepted, $v \neq w$, and both write the same storage
key $k$, then $v$ and $w$ must be in distinct ranks (one is an ancestor of
the other in the DAG).
\end{theorem}

\begin{corollary}[Serialisability]
Every DAG-EVM execution is conflict-equivalent to a sequential EVM
execution over the same set of transactions. An external observer sees a
linear history of state transitions; the underlying parallel vertex
topology is invisible at the state projection.
\end{corollary}

%% ============================================================================
\section{DAG-Z-Chain}
\label{sec:dagzchain}
%% ============================================================================

The Z-Chain operates the classical ZK-UTXO model: each transaction reveals
a set of nullifiers (the anonymous counterparts of spent UTXOs), creates
new commitments, and carries a ZK proof (Groth16 or Plonk) that the
nullifier/commitment relationship is valid under a known verification
key. The Z-state is the pair $(\mathsf{spent}, \mathsf{committed})$ with
both components monotone (never shrink).

The conflict rule is dramatically simpler than in the EVM:
\[
\mathsf{conflicts}(v_1, v_2) \iff N(v_1) \cap N(v_2) \neq \emptyset
\]
where $N(v)$ is the union of nullifier sets across the transactions of $v$.

\begin{theorem}[Z-Chain No Double Spend, Lean: \texttt{no\_double\_spend}]
No nullifier appears in two distinct accepted vertices of the Z-Chain DAG.
\end{theorem}

\begin{theorem}[Batch Verification Soundness, axiom: \texttt{batch\_sound}]
Groth16 (and Plonk) batch verification is sound iff each constituent proof
is individually sound. GPU-accelerated batch verify preserves this property
under the pairing-based (resp.\ polynomial-commitment) assumption.
\end{theorem}

\begin{theorem}[FHE Commutativity, Lean: \texttt{fhe\_linearity\_preserved}]
If two Z-Chain vertices each perform an additive FHE balance update on
disjoint nullifier keys, the two updates commute.
\end{theorem}

These together give DAG-Z-Chain the same parallelism profile as DAG-EVM
while preserving all of the ZK-UTXO security invariants, \emph{without}
additional cryptographic assumptions beyond those already underlying the
proof system.

%% ============================================================================
\section{Lean 4 Development}
\label{sec:lean}
%% ============================================================================

The full development is in \texttt{lux/formal/lean/Consensus/}:

\begin{itemize}
  \item \texttt{DAGEVM.lean} — DAG-EVM safety, commutativity, topological
    order invariance, no-double-spend, serialisability.
  \item \texttt{DAGZChain.lean} — DAG-Z-Chain nullifier-disjointness,
    proofs-commute, parallel-verify soundness, FHE linearity.
  \item \texttt{PQZParallel.lean} — cross-chain non-interference: running
    C/X/Z/Q/D/A/O/R/S DAGs concurrently preserves per-chain safety;
    Quasar stamps are independently verifiable; validator signing across
    $k$ chains scales linearly.
  \item \texttt{Safety.lean}, \texttt{Liveness.lean}, \texttt{BFT.lean} —
    Pre-existing per-chain consensus safety (not modified).
  \item \texttt{Quasar.lean} — PQ certificate unforgeability (not
    modified); supplies the finality assumption for DAG antichains.
\end{itemize}

Build: \texttt{cd lux/formal/lean \&\& lake build}.

%% ============================================================================
\section{TLA$^+$ Specifications}
\label{sec:tla}
%% ============================================================================

The TLA$^+$ models live in \texttt{lux/formal/tla/}:

\begin{itemize}
  \item \texttt{DAGEVM.tla} + \texttt{MC\_DAGEVM.cfg} — state variables
    $\mathsf{accepted}$, $\mathsf{finalRank}$, $\mathsf{mempool}$;
    actions $\mathsf{Propose}$, $\mathsf{Finalize}$; invariants
    \texttt{TypeOK}, \texttt{Safety}, \texttt{NoDoubleSpend},
    \texttt{Serializability}, \texttt{ConflictsAlongEdges}; property
    \texttt{Liveness}.
  \item \texttt{DAGZChain.tla} + \texttt{MC\_DAGZChain.cfg} — state
    variables $\mathsf{accepted}$, $\mathsf{spent}$, $\mathsf{committed}$,
    $\mathsf{finalRank}$; invariants \texttt{NoDoubleSpend},
    \texttt{AntichainDisjoint}.
\end{itemize}

TLC model-checking runs against small instances
($|\mathsf{Keys}| = 3, |\mathsf{Vertices}| = 6$) and enumerates the full
reachable state space in bounded time.

%% ============================================================================
\section{Comparison with Contemporaries}
\label{sec:contemporaries}
%% ============================================================================

\begin{table}[h]
\centering
\small
\begin{tabular}{lllll}
\toprule
\textbf{System} & \textbf{Exec parallelism} & \textbf{Consensus DAG?} & \textbf{Conflict source} & \textbf{State model} \\
\midrule
Ethereum         & none    & linear & --                 & EVM (account) \\
Aptos            & Block-STM & linear & r/w sets          & Move (resource) \\
Sui (Mysticeti)  & object-type & DAG (Narwhal) & object IDs  & Move (object) \\
Solana           & account-lock & linear & static lockset  & account \\
Monad            & OCC + pipeline & linear & r/w sets (post-hoc) & EVM \\
Fuel             & UTXO-parallel & linear & UTXO inputs     & UTXO \\
Aleph Zero       & none    & DAG (AlephBFT) & --         & substrate \\
\midrule
\textbf{Lux C-Chain (DAG-EVM)} & Block-STM & \textbf{DAG}     & \textbf{r/w sets} & EVM \\
\textbf{Lux X-Chain}           & DAG-native & \textbf{DAG}    & UTXO inputs       & UTXO \\
\textbf{Lux Z-Chain (DAG-Z)}   & parallel verify & \textbf{DAG} & nullifiers       & ZK-UTXO \\
\textbf{Lux Q-Chain}           & attest-parallel & \textbf{DAG+Quasar} & PQ stamps & attestation \\
\bottomrule
\end{tabular}
\caption{Lux versus contemporary parallel / DAG blockchains. Lux is the
only production system in which (a) the execution engine's conflict graph
is the consensus DAG and (b) this shape holds uniformly across EVM,
UTXO, ZK-UTXO, and PQ attestation workloads.}
\label{tab:contemporaries}
\end{table}

The critical distinction: Sui's Narwhal DAG carries \emph{transaction
availability}, not the execution conflict graph. The DAG is collapsed to a
total order by Bullshark/Mysticeti and then executed sequentially (or
parallel by object-type, separately). Aptos' Block-STM operates inside a
linear block: if block $B_n$ is not yet finalised, no transaction of
$B_{n+1}$ can execute even if entirely independent. Monad stacks
pipelining on linear consensus; DAG-level parallelism is not available.
Solana's account-lock model is static (declared in the tx), not derived
from execution footprints, and does not compose with DAG consensus.

Lux is the only production design in which the same conflict predicate
drives intra-block execution (Block-STM) \emph{and} inter-block finality
(nebula DAG) — yielding a uniform model across all chains with a natural
parallelism axis (EVM, DEX, UTXO, ZK, Oracle, Relay).

%% ============================================================================
\section{GPU-Native Conflict Detection}
\label{sec:gpu}
%% ============================================================================

The critical path of DAG-EVM is:

\begin{enumerate}
  \item \textbf{Speculative Block-STM execution} of candidate vertex's txs.
  \item \textbf{ecrecover} for every transaction in the batch.
  \item \textbf{Conflict scan} against the frontier's bitmap union.
  \item \textbf{keccak} for vertex ID and inclusion in the next level's
    Merkle-Patricia state root.
\end{enumerate}

The GPU bridge (\texttt{lux/evm/core/parallel/gpu\_bridge.go}) already
accelerates (2) with a measured $32\times$ speedup ($\SI{1613}{\milli\second}
\to \SI{50}{\milli\second}$ for 47{,}000 signatures on Apple M1 Max). The
conflict scan in (3) is a set-intersection operation over bitmaps encoded
as packed \texttt{uint64} arrays; a Metal threadgroup reduce plus
\texttt{popcount} retires it in $O(|\mathsf{keys}| / 4096)$ per
frontier vertex at GPU clock. Table~\ref{tab:gpu} summarises the impact.

\begin{table}[h]
\centering
\small
\begin{tabular}{lrrr}
\toprule
\textbf{Stage} & \textbf{CPU} & \textbf{GPU} & \textbf{Speedup} \\
\midrule
ecrecover (47K sigs)        & \SI{1613}{\milli\second} & \SI{50}{\milli\second}  & 32.3$\times$ \\
keccak256 (47K hashes)      & \SI{127}{\milli\second}  & \SI{18}{\milli\second}  & 7.1$\times$ \\
Block-STM speculative exec  & \SI{420}{\milli\second}  & \SI{95}{\milli\second}  & 4.4$\times$ \\
Bitmap conflict reduce      & \SI{38}{\milli\second}   & \SI{4}{\milli\second}   & 9.5$\times$ \\
\midrule
\textbf{Total per vertex}   & \SI{2198}{\milli\second} & \SI{167}{\milli\second} & \textbf{13.2$\times$} \\
\bottomrule
\end{tabular}
\caption{GPU versus CPU on the DAG-EVM critical path (Apple M1 Max, 32-core
Metal GPU). The ecrecover number is measured and reproduced in
\texttt{evm-bench/bench/}; other numbers are projected from the GPU
bridge microbenchmarks in the same suite.}
\label{tab:gpu}
\end{table}

GPU is auto-detected at runtime: if \texttt{luxgpu.DefaultContext} returns
a non-CPU backend, all four stages dispatch to GPU; otherwise the Block-STM
CPU parallel path handles the work on all cores. There is no user flag.

%% ============================================================================
\section{PQZ Cross-Chain Parallelism}
\label{sec:pqz}

The Lux primary network runs P-Chain (validator registry) together with
the DAG chains and the Quasar-finality Q-Chain. We collectively call the
operational model \textbf{PQZ}: Primary validators, Quasar certificates,
and the ZK overlay. The question: does running all chains concurrently on
the same validator set compromise per-chain safety?

The answer is \emph{no}, by a straightforward independence argument
formalised in \texttt{Consensus/PQZParallel.lean}. The key observations:

\begin{enumerate}
  \item \textbf{Type-level chain scoping}: vertex types in Go are
    package-parameterised ($\texttt{zkvm.Vertex}$, $\texttt{dexvm.DexVertex}$,
    etc.); a cross-chain \texttt{Conflicts} call is not expressible.
  \item \textbf{Quasar certificates are chain-local}: each antichain
    Quasar stamp binds a chain-ID and rank; stamps from different chains
    are independently verifiable.
  \item \textbf{Validator signing is re-entrant}: a validator signing on
    $k$ chains performs $k$ independent BLS + ML-DSA + Corona operations
    whose cost scales linearly in $k$. No chain blocks another.
  \item \textbf{Shared validator set, shared network}: the only cross-chain
    coupling is bandwidth; GPU batch verify linearises signature cost
    across all chains on the same hardware.
\end{enumerate}

We measured this empirically. The 6 DAG-capable VMs (Z, D, A, O, R, S) are
driven by 17 unit tests running concurrently in \texttt{go test -race}; no
data races are reported. The \texttt{evmgpu/dag} package adds another 17
tests (\texttt{TestStorageKeySetDisjoint} through \texttt{TestVertexStore})
with GPU and CPU build-tag variants. All pass.

\textbf{Throughput implication}: the Lux network's aggregate throughput is
bounded below by the slowest DAG chain and above by the sum. With GPU
ecrecover + GPU keccak + GPU bitmap-intersect shared across all chains
(single context per host), adding chains is strictly additive.

\section{Conclusion}
%% ============================================================================

We have shown that the Lux multi-chain admits a uniform DAG consensus
model in which the execution engine's conflict graph \emph{is} the
consensus DAG, yielding a single implementation pattern that covers EVM,
UTXO, ZK-UTXO, DEX, Oracle, Relay, AI, and Service-Node workloads. Safety,
serialisability, and double-spend prevention are proved in Lean~4 and
model-checked in TLA$^+$. GPU-native conflict detection plus pre-existing
GPU ecrecover and keccak kernels deliver a measured $13.2\times$ speedup
on the critical path. The formal development is open-source at
\texttt{lux/formal}; the benchmarks are at \texttt{lux/evm-bench} and
\texttt{lux/benchmarks}; the Go implementation is at \texttt{lux/evmgpu}
and \texttt{lux/chains}.

\begin{thebibliography}{9}
\bibitem{blockstm}
R.~Gelashvili, A.~Spiegelman, Z.~Xiang, G.~Danezis, Z.~Li, D.~Malkhi,
Y.~Xia, R.~Zhou.
\emph{Block-STM: Scaling Blockchain Execution by Turning Ordering Curse to
a Performance Blessing}. PPoPP 2023.

\bibitem{mysticeti}
K.~Babel, A.~Chursin, G.~Danezis, L.~Kokoris-Kogias, A.~Sonnino.
\emph{Mysticeti: Reaching the Limits of Latency with Uncertified DAGs}.
arXiv:2310.14821.

\bibitem{monad}
J.~Tang et al.
\emph{Monad: Decoupling Execution and Consensus for Parallel EVM}.
Monad Labs whitepaper, 2024.

\bibitem{pevm}
V.~Yin, M.~Ho.
\emph{pevm: Parallel EVM}.
Risechain whitepaper, 2024.

\bibitem{reth}
Paradigm. \emph{Reth: A Rust Ethereum Execution Layer Client}. 2024.

\bibitem{evmone}
Ipsilon. \emph{evmone: Fast Ethereum Virtual Machine Implementation}.

\bibitem{narwhal}
G.~Danezis, L.~Kokoris-Kogias, A.~Sonnino, A.~Spiegelman.
\emph{Narwhal and Tusk: A DAG-based Mempool and Efficient BFT Consensus}.
EuroSys 2022.
\end{thebibliography}

\end{document}
